% PROPOSAL ONLY. This file is not included by the manuscript.
% Replace the score definition in Section 5 and reconcile the separate
% profiling/5QI paragraph, metric names, captions and claims listed in TASK2G_REPORT.md.

The profiling target is a study-specific composite quality score, rather than
an indicator of formal 5QI compliance. Historical score labels were reconstructed
from the surviving scoring notebook and the measurement tables. For profiling
observation $i$, let $a_i$ be the recorded per-direction offered-load objective
and let $r_i^{
\mathrm{DN,rx}}$ be the stored DN-side receive-throughput counter, both in Mbps.
Let $d_i$ denote \texttt{downlink\_delay\_ms}, $j_i$ denote
\texttt{RAN\_jitter} in microseconds, and $\ell_i$ denote the stored fraction
\texttt{RAN\_PLR\_interval}. The receive-counter direction and the loss-denominator
construction have not been authenticated end to end; this definition preserves
the recorded measurement fields without relabelling them as an exclusively
downlink service measurement.

Define $\operatorname{clip}(x;a,b)=\min\{b,\max\{a,x\}\}$. The component scores are
\begin{align}
 S_i^{\mathrm{thr}} & = \operatorname{clip}\!\left(
       \frac{r_i^{\mathrm{DN,rx}}}{a_i+\varepsilon};0,1\right),
       &&\varepsilon=10^{-9}\ \mathrm{Mbps}, \nonumber\\
 S_i^{\mathrm{lat}} & =1-\frac{\operatorname{clip}(d_i;0,40)}{40},
       &&d_i\text{ in ms}, \nonumber\\
 S_i^{\mathrm{jit}} & =1-\frac{\operatorname{clip}(j_i;0,10)}{10},
       &&j_i\text{ in }\mu\mathrm{s}, \nonumber\\
 S_i^{\mathrm{loss}} & =1-\frac{\operatorname{clip}(100\ell_i;0,5)}{5},
       &&100\ell_i\text{ in percent}.
 \label{eq:qos_components_recovered}
\end{align}
The historical weighted score is
\begin{equation}
 Q_i=0.40S_i^{\mathrm{thr}}+0.20S_i^{\mathrm{lat}}
       +0.20S_i^{\mathrm{jit}}+0.20S_i^{\mathrm{loss}}.
 \label{eq:qos_weighted}
\end{equation}
The numerical values 40~ms, 10~$\mu$s and 5\% are score-normalization caps used
in the experiments; they are not asserted to be standardized service budgets.
The previously stated three-component weights and 5QI-derived interpretation
were a methodological reporting discrepancy. The experimental labels use the
four-component definition above and are retained unchanged.

The twin's QoS regressors approximate this composite score from per-instance
offered load. Thus, for an executed configuration with mode $m_t$ and $k_t$
balanced instances, the quantity used by the controller is the model prediction
$\widehat Q_t=q_{m_t}(\rho_t/k_t)$. We report the frequency and severity of
\emph{predicted composite-score shortfalls} below $\tau=0.90$; separate profiling
validation compares predictions with the measured score labels. The reward uses
a soft penalty on these modeled shortfalls. Neither a prediction above $\tau$
nor the reward formulation establishes an end-to-end delay, packet-error-rate,
session-continuity or formal 5QI guarantee. The exploratory PCA-derived score
\texttt{QoS\_dd} is not the target of the QoS profiling models.

% Separate replacement for the profiling-direction/5QI paragraph:
% The profiling exports contain DN- and RAN-side transmit/receive counters.
% The recorded per-direction offered objective is used as the model input.
% The study evaluates the recovered composite score defined above, rather than
% a standardized 5QI conformance test. Device/flow attribution and loss-counter
% construction are limitations; do not describe the whole campaign as downlink-only.
% Standardized QoS characteristics can remain background context with
% \cite{3gpp23501}, but the score's caps are not derived from that table.

% Grouped-protocol paragraph: use only after the planned experiments have run.
% For each UPF mode, we use complete sweeps 1--3 for parameter learning,
% sweep 4 for checkpoint and model selection, and sweep 5 for internal testing.
% Input normalizers are fitted on training sweeps only. The internal-test
% results do not select models or extend the candidate search. These exposed
% same-campaign measurements do not constitute a fresh physical holdout.
